\chapter[Buddhism Changes Across Asia]{How Buddhism Kept Changing Across Asia}
\section{A Buddha Whose Laughter Seems Un-Buddhalike}
First, pick up an object---or, strictly speaking, look up at one.

If you have visited any reasonably large Buddhist temple in China, the first hall beyond the entrance probably presented you with a figure like this: an unabashedly fat man, his robe open and round belly exposed, laughing so broadly that his eyes narrow to two slits. A couplet often hangs on the pillars beside him, conveying something like "A great belly accommodates all; an open mouth always laughs." Visitors like him, children dare to touch his stomach, and people dropping money into the donation box sometimes nod to him as if greeting a good-natured neighbor.

People at the temple will tell you: this is Maitreya Buddha.

But wait. Open an album of ancient Indian Buddhist art and find an image bearing the same name, Maitreya. You will see someone entirely different: a well-proportioned figure with a quiet expression, standing or sitting, hands held gracefully, displaying a trained and restrained compassion. He seems to have no family resemblance whatsoever to the laughing fat man in the temple. In fact, the Chinese figure is modeled on a wandering monk from Fenghua in Zhejiang during the Five Dynasties period, about eleven hundred years ago. He traveled begging with a cloth sack, smiling at everyone, and left a verse before his death. People later believed him to be an incarnation of Maitreya. His likeness laughed its way into the Hall of the Heavenly Kings, completely replacing the solemn Indian image.

Look farther afield. Japanese Buddhist images have Japanese faces; Thai images have Thai faces. Place the ornate, densely arranged images of Tibetan temples beside Sri Lanka's white stupas and it is hard to imagine a common source. Yet they do share one: around twenty-five hundred years ago, a practitioner in northern India named Shakyamuni and the tradition he began.

This is our chapter's subject. During two thousand years of travel from its homeland across Asia, Buddhism did not transport itself unchanged, inch by inch. It kept changing as it moved: its languages, the faces of its sacred images, its methods of practice, and its dealings with deities. So much changed that an ancient Indian monk brought back to life might stand bewildered at the entrance to a temple in today's Tokyo or Lhasa. Our question is: did these transformations distort Buddhism, or keep it alive?

\section{A Translation Workshop in Chang'an}
Now come with me to a particular scene.

Time: the mid-seventh century CE, during the Tang dynasty. Place: the scripture translation institute at the Great Ci'en Temple in Chang'an. The crown prince ordered this temple built in memory of his mother, and it has ample worshippers and funding. Today, however, its busiest place is not the main hall but a large room: the translation workshop.

A thin-faced monk sits at its center. His name is Xuanzang. Twenty years earlier, as a young monk, he had done something bordering on illegal. In the early Tang, ordinary people were forbidden to leave the frontier without permission. He slipped across the border among refugees fleeing famine, then traveled alone through deserts, snowy mountains, and the Gobi, covering thousands of kilometers to study Buddhism in India. After years of study at Nalanda, India's most famous Buddhist center of learning, he made the journey back with hundreds of Sanskrit scriptures. The emperor pardoned him and gave him a task: translate them into Chinese.

What lies before Xuanzang now is not merely a book but an entire production line. Historical accounts describe an astonishingly detailed division of labor in Tang translation workshops. One person held the Sanskrit original and read it sentence by sentence; another checked pronunciation; another wrote an initial Chinese draft from the oral translation; others rearranged and polished its sentences. Someone checked its consistency with doctrine, and someone else neatly copied the final version. A major scripture could occupy dozens or even hundreds of people for years. Imagine the sounds in that room: Sanskrit recitation, brushes rustling over paper, and occasional eruptions of argument as two learned monks grew red-faced over the translation of a single word.

What were they arguing about? Take an example you may know. The bodhisattva who rescues people from suffering is popularly called Guanshiyin in China. This translation of the name had circulated for centuries before Xuanzang. But after checking the Sanskrit, Xuanzang concluded that earlier translators had been mistaken: the original meaning called for Guanzizai, suggesting freedom gained through inward contemplation. He insisted on Guanzizai in his own translations. What happened? Two hundred years, then a thousand years passed, and ordinary people kept burning incense, bowing, and invoking "Guanshiyin Bodhisattva," disregarding the great scholar's research. The old name had taken root in millions of people's everyday language and could not be pulled out.

Remember this argumentative workshop. Half of Buddhism's Asian journey took place on the road; the other half was argued into being, word by word, in rooms like this.

\section{Is Buddhism Outside India Still Buddhism?}
Let us first set out the tension, without rushing to an answer.

Buddhism originated in India. According to tradition, before renouncing worldly life Shakyamuni was a prince of the Shakya clan who abandoned palace life to practice and eventually awakened beneath the bodhi tree. These are matters of belief recorded in the scriptures. What historians can establish is that around the fifth century BCE, an influential shramana teacher and the community he founded appeared in the Ganges basin. Over the following centuries, Buddhism grew from a local community of practitioners into a major tradition across the subcontinent, divided into numerous schools, and spread abroad: south to Sri Lanka and Southeast Asia; north through Central Asia into China; from China or by sea to the Korean peninsula and Japan; and into the Tibetan Plateau from both India and the Chinese heartland.

A pointed question therefore arose---one that Buddhist scholars themselves were asking at the time:

\textbf{After traveling so far from India and changing so much, was Buddhism still Buddhism?}

Two positions lie inside this question. One says: of course. The Buddha's teachings are like seeds. In different soils they grow into different trees; however much the trees vary in appearance, what was in the seed remains unchanged. The other says: not necessarily. If a religion's language, scriptures, rituals, and sacred images have all changed, even its answer to "what counts as practice," perhaps all that connects it to its source is a name.

More troublingly, the second position has an "official version." Buddhism contains a widespread prophetic account called the "age of Dharma decline": after the Buddha's passing, his teaching will gradually fade, like the setting sun. Later ages grow darker, explanations become mistaken, and practice becomes mere form. On this account, all later developments are inherently suspect. New does not necessarily mean wrong, but the thought "the farther from the Buddha, the more suspicious" is written into the tradition itself. Buddhism, then, is not merely a religion that changes continually; it is a religion \textbf{anxious about change}.

An outside observer might ask the opposite question: what in the world does not change? Languages change, states change, even people's food changes. Why should a tradition alive for twenty-five hundred years remain exactly the same? A tree that grew no new branches for two thousand years would have died long ago.

\subsection{First Meet Four Pieces of Baggage}
To discuss whether something changed, we must first know what it carried when it set out. Early Buddhism's basic equipment consisted roughly of four things, each repeatedly refitted during its later travels.

\textbf{The sangha}: a self-governing community of renunciant practitioners. It resembles both a lifelong boarding school and a modern open-source community. There is no king-like supreme leader; major matters are discussed in assemblies. Members give up family life and private property and live on supporters' offerings. After the Buddha died, this community, rather than any one person, preserved and continued his teaching.

\textbf{Monastic discipline}: the sangha's internal rules, ranging from fundamental prohibitions on killing and stealing to details about storing robes and bowls or where to stay during the rainy season, numbering in the hundreds or thousands. Do not think of them as a list of restraints. They are the community's operating system: an organization without shared rules cannot survive two generations. Discipline is one of the technical secrets behind Buddhism's survival for over two thousand years. Notice too that these rules mainly govern monastics. Lay followers have lighter versions such as the Five Precepts. Buddhism was therefore a "two-track system" from the outset, with different doors for intensive and ordinary participants.

\textbf{Meditation}: methods for systematically training attention and the mind. By analogy, listening to doctrinal instruction is like reading an exercise manual; meditation is actually entering the gym and practicing set after set, training the muscle that clearly observes thoughts arising and passing. Meditation is neither absent-minded staring nor blanking out. It is a highly alert technique of inward observation---a point we will revisit near the chapter's end.

\textbf{Scriptures}: Buddhist literature is called the Tripitaka, or "Three Baskets": sutras, containing teachings; vinaya, containing disciplinary rules; and abhidharma, containing subsequent analysis and discussion. It is one of the world's largest bodies of religious literature and has never settled into a single universally authoritative collection. Different languages and schools preserve canons with differing contents and wording. Remember this when we read the Platform Sutra later.

The four pieces of baggage are packed. Now the question is: when they travel from the Ganges basin to Chang'an, Kyoto, Lhasa, and Bangkok, across deserts, seas, and entirely different civilizations, which will arrive unchanged?

Before continuing, take a position of your own. Is the smiling, potbellied Maitreya in the Hall of the Heavenly Kings a sign of Buddhism's "corruption" or its "ingenuity"?

\section{Who Protects It, Who Drives It Out?}
To understand how Buddhism changes, first consider the people whose competing pulls bring about change. People in different positions see entirely different things in the same event.

\textbf{Voice one: the king's calculations.}

In the third century BCE, Ashoka of India's Maurya dynasty marked a turning point. As a young ruler he waged wars of conquest. According to inscriptions he himself ordered carved into rock, one campaign's mass deaths, injuries, and displacement filled him with remorse. He subsequently supported Buddhism, sent people throughout his lands and abroad to spread the teaching, and had edicts carved on stone pillars and rock faces. One passage conveys this: of all victories, only victory won through dharma is true victory (Ashoka's Major Rock Edict XIII). Notice what happened: an emperor took a religion's central term, dharma---in Buddhist usage, both the Buddha's teaching and right conduct or order---and made it a banner of government. Buddhism became not only the sangha's undertaking but an imperial resource. This was protection and transformation together. From Ashoka onward, almost every major Buddhist expansion across Asia had a political authority behind it.

\textbf{Voice two: the scholar-official's anger.}

Han Yu of Tang China stood at the opposite pole. In the early ninth century CE, the emperor planned to bring what was said to be a bone from the Buddha's finger into the palace for worship. The country was swept with excitement; some people burned their fingers or scorched their scalps to display devotion. Han Yu submitted a fiercely opposing memorial, beginning:

\begin{quote}
"I respectfully submit that Buddhism is merely a teaching of foreign peoples." (Han Yu, Memorial on the Buddha's Bone)
\end{quote}
His point was that Buddhism was a foreign teaching. His further objections included these: monks neither plow nor weave, consuming society's food without producing grain; renunciation abandons the family and violates the principles of loyalty and filial duty. The furious emperor demoted him to distant Chaozhou.

This episode is often told as "ignorant superstition versus reason," but try stating Han Yu's case fully. It is a worthwhile exercise. His accounts were concrete: after the High Tang, large amounts of land and labor flowed into monasteries, which paid no taxes, genuinely straining state finances. Nor were his ethical concerns empty words. In a society built around family obligations, an institution requiring people to leave their parents, remain unmarried, and have no children was a structural affront. Decades later, Emperor Wuzong launched a major suppression of Buddhism, the mid-ninth-century Huichang Persecution, destroying tens of thousands of temples and forcing hundreds of thousands of monks and nuns back into lay life. He used precisely the kind of accounting Han Yu had offered. You may disagree with Han Yu, but his opposition was not superstition: it was a complete view of how the state should work.

\textbf{Voice three: the gods' territory.}

Now Japan. Buddhism arrived in the Japanese islands in the mid-sixth century CE through Baekje on the Korean peninsula, immediately dividing the court. The Soga clan advocated acceptance as advanced culture; the Mononobe and Nakatomi clans strongly opposed it. Their reasoning was that after centuries of worshipping local deities, suddenly turning to foreign gods might provoke the native gods' anger. Traditional accounts say that Buddhist images were first enshrined and then an epidemic broke out. Opponents immediately declared, "See, the gods are angry," and the images were thrown into a river. When the epidemic continued, they were retrieved. After decades of struggle, Buddhism prevailed, but how it prevailed matters: it never replaced the native deities, instead learning to \textbf{coexist} with them. We will examine this more closely in the "Further Challenge" section.

\textbf{Voice four: practitioners at a fork in the road.}

Within the tradition, disagreements were just as deep. Faced with the same question---"What does practice ultimately depend on?"---different traditions offered almost opposite answers:

\begin{itemize}
\item \textbf{Chan}, or Zen, says: yourself, and not your reading. It proclaims "a special transmission outside the teachings, not founded on written words," locating awakening beyond scripture, in sitting meditation, work, or even a master's unexpected blow or shout. Bodhidharma, who came from India to China around the fifth century, is honored as its first patriarch. But the person who truly made Chan "Chinese" was the Tang figure Huineng, said to have been an illiterate southern woodcutter.
\item \textbf{Pure Land} says: yourself? Do not overestimate your strength. In the age of Dharma decline, liberation through one's own power is like rowing upstream. Faithfully reciting "Homage to Amitabha Buddha" and relying on Amitabha's vows for rebirth in the Land of Bliss is the reliable path. Recitation moved practice from remote mountain monasteries to field banks and kitchen stoves. People unable to read or spare time for seated meditation received an equal admission ticket for the first time.
\item \textbf{Esoteric Buddhism}, or Vajrayana, says: both paths are too slow. Using rituals, hand gestures, mantras, and visualization, it advocates attainment in this very life and body, without waiting for gradual progress. It flourished in Tang Chang'an and later rooted deeply in Japan, through the esoteric tradition Kukai brought back, and on the Tibetan Plateau. In Tibet it combined with local traditions and developed institutions unknown elsewhere, such as reincarnating lamas.
\end{itemize}
All three voices have a case, and each finds fault with the others. Pure Land followers regard Chan as a game for a gifted few; Chan masters regard recitation as intellectual laziness; esoteric teachers think both are taking unnecessary detours. Remember: Buddhism has never spoken with one voice. Its internal disagreements are no smaller than its arguments with outsiders.

\section{Thinking Bridge: Travel with a Word}
When asking whether a tradition changed, abstract assertions that it did or did not are unhelpful. Here is a portable tool: \textbf{tracing a word's travels}. Choose an everyday word, find out whether it came through translation, and ask four questions. Where was it originally? Who brought it here? What was added along the way? What did it become after arrival?

Try a Chinese word used daily: shijie, "world." Today it means the entire globe or human society. But it comes from Buddhist translation: shi refers to the flow of time---past, present, and future---while jie refers to spatial directions---east, south, west, north, above, and below. Translating monks compressed a whole Buddhist conception of time and space into two Chinese characters. Whenever Chinese speakers say "World Cup" or "worldview," they unwittingly use a translation tool more than a thousand years old.

Try another: chana. Today chana jian means "in an instant." It comes from a Sanskrit transliteration, originally an extremely short unit of time in Indian traditions. Then there is fannao, "vexation": in Buddhist scriptures it is a technical term for the greed, hatred, and delusion that disturb the mind. When a Chinese speaker says, "I'm so fed up," they use its everyday descendant. Fangbian, "convenience"; juewu, "awakening"; pingdeng, "equality"; jietuo, "release"---a great many Buddhist translation terms have become so thoroughly part of Chinese that nobody remembers they were once newcomers.

Now see the tool's real power by asking what was added in transit. During the Wei and Jin periods, as Buddhist scriptures first arrived in large numbers, translators found no ready Chinese equivalent for the Sanskrit concept of "emptiness," shunyata, and borrowed the Daoist wu, "nonbeing." For a Sanskrit term denoting the path to awakening, they borrowed dao, "the Way." Using old Chinese concepts as translation crutches later became known as geyi, "matching concepts." The crutches helped but had side effects. For a time, Wei-Jin readers genuinely treated Buddhism as a relative of their own metaphysical learning, assuming the Buddha and Laozi were saying the same thing. Translators such as Kumarajiva, the Kuchean scholar who directed translation in Chang'an in the late fourth and early fifth centuries, and Xuanzang gradually removed these crutches and substituted newly coined terms. Yet notice: even correcting mistakes was an act of creating new Chinese.

This tool is not limited to Buddhism. The Chinese terms for "science," "democracy," "economy," and "society" underwent similar journeys when borrowed back through Japanese in modern times. "Involution," traveling from an anthropological work into today's schoolchildren's speech, has also turned several semantic somersaults. The next time someone argues about a word's "original meaning," you can ask: which stop on its journey do you mean?

\section[Reading the Heart Sutra]{A Storm in 260 Characters: Reading the Heart Sutra}
Now let us read a short primary source. The Buddhist canon is vast; we will choose one of its shortest and most widespread texts, Xuanzang's translation of the Prajnaparamita Heart Sutra. At around 260 Chinese characters, it is nevertheless among the most frequently recited scriptures in temples from Xi'an to Tokyo. It begins:

\begin{quote}
When Guanzizai Bodhisattva practiced the profound perfection of wisdom, he clearly saw that all five aggregates were empty and passed beyond all suffering and distress. Shariputra, form does not differ from emptiness, nor emptiness from form; form is emptiness, and emptiness is form. (Prajnaparamita Heart Sutra, translated into Chinese by Xuanzang during the Tang dynasty)
\end{quote}
First let us unpack the literal terms. The "five aggregates" are Buddhism's analysis of a person: form, meaning body and matter; feeling; perception; volitional formations; and consciousness. Bundled together, they make what we take to be "me." The "emptiness" in "form is emptiness" does not mean that nothing exists. It means that nothing possesses an independent, unchanging nature: everything arises under conditions and passes away as conditions change. The argument, broadly, is that if even the five constituents of "me" are empty, the suffering produced by attachment to "me" loses its foundation.

Notice a detail that encapsulates this chapter. Xuanzang writes "Guanzizai" here. Remember the argument in the Chang'an workshop? He insisted on this name to correct the older Guanshiyin. A curious situation resulted: Chinese people recited Xuanzang's prized Heart Sutra while continuing to call its opening figure by the name he rejected. The translator won the scripture; ordinary people won the name. Authority and habit each won half.

Now a more dramatic source. The Platform Sutra, the only work by a Chinese monk honored with the title "sutra," tells a story about choosing a successor. The fifth patriarch, Hongren, asked his disciples to compose verses. His senior disciple Shenxiu wrote:

\begin{quote}
The body is the bodhi tree; the mind is like a bright mirror's stand. Diligently wipe it at every moment; let no dust settle upon it.
\end{quote}
The point is that practice resembles cleaning a mirror: polish it diligently every day so dust cannot settle. Hearing this, Huineng, a laborer who pounded rice, asked someone to write another verse for him. In the commonly circulated Platform Sutra it reads:

\begin{quote}
Bodhi originally has no tree; the bright mirror has no stand. From the beginning there is not a single thing; where could dust settle?
\end{quote}
The point is that if everything is empty, where are the tree, the stand, or the dust to wipe away? Hongren passed his robe and bowl to the illiterate Huineng. A more important detail follows. Several versions of the Platform Sutra survive, and in the earliest Dunhuang manuscript Huineng's verse differs from the familiar version. Even this verse rebutting Shenxiu was rewritten during later copying and transmission. A story about the unreliability of words was itself transmitted through an unreliable text. This is not irony but evidence: the transmission of scripture is itself continuing re-creation.

\section{Why It Changed: Three Explanations}
We now have enough pieces to lay out genuinely different explanations. First specify the facts to be explained: Buddhism arose in India in the fifth century BCE and flourished abroad, yet by around the twelfth and thirteenth centuries it had largely disappeared from India itself, its communities dispersed and great monasteries abandoned. The same tree withered in its homeland and became forests abroad, each foreign forest looking different. Why?

\textbf{Explanation one: dilution---the farther from the source, the greater the distortion.}

On this account, transmission is a game of whispered messages. Every translation and cultural crossing loses another layer of information; Buddhism reaching Japan is a diluted version of the original concentrate. This explains several things, including the persistence of Dharma-decline thinking and eminent monks' recurring calls to return to the Buddha's original intention. But its limitation is fatal: it assumes a fixed "original concentrate." Buddhism began dividing into schools soon after the Buddha's death, and their records disagree even on what he said. Without a fixed original, dilution has no measuring standard.

\textbf{Explanation two: adaptation---change is not loss but a way of surviving.}

On this account, religions resemble species: entering a new environment requires change or death. Encountering a Chinese society centered on family and filial duty, Buddhism translated scriptures about parents' kindness and developed ceremonies to aid ancestors after death. Encountering the Tibetan Plateau's local beliefs, it absorbed protective deities and reincarnating lamas. Encountering Shinto in Japan, it learned to share temples with native gods. Changing shape demonstrates vitality. Its decline in India can partly be explained by other traditions occupying its ecological niche there. This explanation is stronger than dilution, but it has a blind spot: it makes religion too much like an organism, adapting and evolving automatically, and obscures \textbf{the people making decisions}---Xuanzang translating scriptures, emperors suppressing Buddhism, ministers throwing images into rivers and retrieving them. Adaptation does not happen automatically; it consists of particular struggles and negotiations.

\textbf{Explanation three: institutions---Buddhism's survival depends on its relationships with organizations, money, and power.}

This explanation shifts position, asking not how faithfully doctrine is preserved but what keeps the sangha alive. Buddhism's skeleton is its monastic community. Discipline, with rules from prohibiting killing to distributing food and robes, maintains internal order; donations sustain its livelihood. This structure is light and flexible, well suited to travel, but has a vulnerability: heavy dependence on supporters. In India, great monasteries such as Nalanda depended on royalty and wealthy merchants. After invading armies destroyed them in the late twelfth century, communities lost protection and lacked roots in village life, so Buddhism rapidly collapsed. In Japan, by contrast, temples were deeply tied to political authorities and landed estates. In Southeast Asia, they became embedded in every village, with boys' temporary ordination a universal rite of passage. The sangha's roots entered society's very muscles and could scarcely be pulled out. This sharp explanation shows why the Huichang Persecution could badly damage Chinese Buddhism and why Thai Buddhism is so difficult to dislodge. Its limitation is that reducing everything to economics and politics can obscure the real force of faith. Financial statements cannot fully explain people burning their fingers as offerings or traveling thousands of kilometers for scriptures.

Each explanation holds part of the truth. Dilution reminds us that transmission has direction and anxiety; adaptation that change is normal rather than accidental; institutions that a spiritual tradition's life and death are written in account books.

\section[Further Challenge: Syncretism]{Further Challenge: Syncretism Is the Norm, Not a Mistake}
Now we introduce the chapter's weightiest concept. Scholars of religion call the interpenetration and recombination of traditions after they meet \textbf{syncretism}. The word long carried a negative sense, suggesting impurity or a makeshift mixture. Research since the twentieth century has reconsidered that judgment: syncretism is not religion's degeneration but its normal condition. A living tradition without it is almost as hard to find as a living language without borrowed words.

Buddhism's Asian history is a gallery of syncretism. In Japan, Buddhism and native Shinto worked out a sophisticated theory called honji suijaku, "original ground and manifested traces." Buddhas were the fundamental reality, and Japanese deities their local traces or manifestations. A mountain could therefore contain both a Shinto shrine and a Buddhist temple; a deity's name could carry a Buddhist honorific. Ordinary people worshipped both kami and buddhas throughout life, marrying at shrines and inviting monks to conduct funerals, without sensing any contradiction. This fusion of kami and buddhas operated for more than a thousand years.

Now consider \textbf{a counterexample that is easy to misread}. Hearing that syncretism is normal, you might conclude that it grows naturally, remains harmonious and stable, and cannot be dismantled. Consider what happened: in 1868, Japan's Meiji government ordered the separation of kami and buddhas. Within a few years, shrines that had housed joint worship for a millennium were "cleansed," Buddhist images removed or smashed, and countless temples abolished---the movement known as haibutsu kishaku, "abolish Buddhism and destroy Shakyamuni." Yesterday's self-evident mixture became today's error demanding correction. This shows two things: syncretism is never simply natural, but always arranged by people; and separation, like combination, is a political decision. The apparently traditional arrangement of distinct Shinto shrines and Buddhist temples in Japan today is actually only a little over a hundred and fifty years old, made by human hands.

Now the Tibetan Plateau. After Buddhism arrived, it interacted over long periods with local Bon traditions, each absorbing the other's rituals and deities. Tibet also developed an institution found nowhere else: reincarnating lamas. After an eminent lama died, people searched for his reincarnated child, transmitting authority across generations in this way. In Myanmar, Buddhism coexists with kingship and popular nat-spirit worship. In Thailand, many men enter monastic life briefly at least once, making the sangha something like a school for the whole population. Each combination is a new arrangement jointly authored by Buddhism and local society, not a photocopy of an original but a shared reorganization.

Syncretism's sharpest critique cuts into the illusion of a pure source. What scholars today call "early Buddhism" is itself \textbf{reconstructed} by later generations from texts such as the Pali canon, which tradition says was first written down in Sri Lanka around the first century BCE, centuries after the Buddha's death. Even the "source," then, is a later work of translation and reconstruction. This does not make truth and falsehood irrelevant. It means that in a living tradition, "returning to the source" is itself a proposition continually rewritten within that tradition. Understand this and every claim to exclusive authenticity will meet an additional measuring stick in your mind.

\section[Mindfulness and Digital Rituals]{Mindfulness, Electronic Wooden Fish, and "Buddha-Like" Living}
This story, more than two thousand years old, is continuing on your phone right now.

Buddhism kept changing across Asia. In the twentieth century it crossed the Pacific and began another journey, taking forms that would probably surprise Xuanzang even more than the potbellied Maitreya.

Consider mindfulness. In the late 1970s, an American doctor named Jon Kabat-Zinn extracted certain Buddhist meditation exercises, removed their religious framework, and made them into an eight-week course to help patients cope with chronic pain and stress. This became the globally influential program of mindfulness-based stress reduction. Mindfulness now appears in hospitals, schools, armies, and Silicon Valley companies. Countless people who never identify as Buddhists close their eyes each day and observe their breathing. The term translates the Pali sati, a component of Buddhism's Noble Eightfold Path. Yet these classes teach neither the Four Noble Truths nor rebirth and nirvana; they train attention alone. It is a relay spanning twenty-five hundred years: Indian meditation, repeatedly transformed in China, Korea, and Japan, rewritten as a psychological technique by an American doctor, and finally returned to Asia in workshops beside urban office towers.

Now look around you. Electronic wooden-fish apps are popular among young people: tap the screen and "merit increases by one." The Chinese expression foxi, "Buddha-like," became an internet favorite a few years ago, describing a relaxed, noncompetitive willingness to let things take their course. This expression has traveled too. Borrowed from a Japanese media label for "Buddha-like men," it changed meaning again on reaching China. There is also AI: people use large language models to interpret Buddhist scriptures, and some temples use robots to explain the Dharma.

The question from our opening now returns in new clothes. Is mindfulness without the Four Noble Truths or rebirth still Buddhism? Different people give sharply opposed answers. A psychological clinician may say: whether it is Buddhism does not matter if it helps someone sleep. Some monks warn that cutting practice away from liberation from birth and death and turning it into stress relief is the latest form of Dharma decline. Other Buddhist scholars remain calmer: when has Buddhism not changed? In India it involved philosophical debate; in China, rites for ancestors; in Japan, funeral services; in America, psychological techniques. Every transformation provoked anguish, and every time it survived.

Even concern about impurity has two faces. One side says mindfulness introduces hundreds of millions to awareness training, an immeasurable benefit. The door is open; those wishing to go deeper can enter. The other says that translating "freedom from suffering" into "stress reduction" quietly replaces more than terminology: it replaces an entire answer to why humans suffer. Buddhism says suffering is rooted in attachment; the stress-reduction industry says you suffer because you have not bought enough of its courses. Neither side is foolish.

\section[Your Question: The Changing Ship]{Your Question: A Ship Whose Planks Keep Changing}
Finally, fit all this chapter's pieces into an ancient thought experiment.

The Greeks told the puzzle of the Ship of Theseus. A ship sails, replacing each plank as it wears out. Years later, none of the original planks remains. Is it still the same ship?

Now replace the ship with Buddhism. Over twenty-five hundred years it has changed languages, from Pali and Sanskrit to Chinese, Tibetan, Japanese, and English; changed the faces of its images, from the restrained compassion of Indian Maitreya to China's laughing Budai; changed its answers about practice, from self-power to other-power to Buddhahood in this very body; changed scriptural wording, even Huineng's verse; and changed neighbors, from Brahmanism to Daoism, Shinto, Bon, and psychotherapy. Not one plank has remained untouched from beginning to end.

Is it, then, still the same ship?

Give your answer and your reasons. Before deciding, weigh these considerations again:

If you say yes, what grounds its sameness? Certain unchanging scriptural sentences, ordination lineages passed down through the sangha, believers' self-identification, or an elusive core called the Dharma? Whichever you choose must face the question: has that not changed too?

If you say no, at which stop did it cease to be Buddhism? When Huineng rewrote Indian meditation? When Pure Land said reciting the Buddha's name was enough? When kami and buddhas shared a hall? When mindfulness entered hospitals? Wherever you draw the line, history contains sincere people living on its other side, convinced that they embody the true tradition.

One smaller measuring stick remains. If you are not a Buddhist, why assume that your criterion for what counts as Buddhism carries more weight than that of an elderly woman who has recited "Homage to Amitabha Buddha" all her life? Or, conversely, why assume hers must prevail?

There is no standard answer. But the next time you hear an argument in any field about what is truly authentic and what no longer is---a dish, a dialect, or a body of thought---you will recognize a negotiation that has already lasted twenty-five hundred years and is destined never to end.
